%Chargement des paquets
\usepackage{amsmath}
\usepackage{amsfonts}
\usepackage{amssymb}
\usepackage{enumerate}
\usepackage{mathtools}
\usepackage{bbm}
\usepackage{xparse, etoolbox}
\usepackage{enumerate}
\usepackage{enumitem}
\usepackage{mathabx}
\usepackage{babel}[french]

%environment
\usepackage{amsthm}
\usepackage{keytheorems}
\usepackage{hyperref} 

%Interface théorème
\newkeytheorem{lem}[
	name = Lemme
]

\newkeytheorem{prop}[
	name = Proposition
]

\newkeytheorem{thm}[
	name = Théorème
]

\newkeytheorem{coro}[
	name = Corollaire
]

\renewcommand*{\proofname}{Démonstration}

\theoremstyle{definition}
\newtheorem{defn}{Définition}

\theoremstyle{definition}
\newtheorem*{nota}{Notation}

\theoremstyle{definition}
\newtheorem*{limi}{Liminaire}

\theoremstyle{remark}
\newtheorem{rmq}{Remarque}

\theoremstyle{plain}
\newtheorem*{fait}{Fait}


%Ensembles de nombres
\newcommand{\N} {\mathbb{N}}
\newcommand{\Ne}{\N^\ast}
\newcommand{\Z} {\mathbb{Z}}
\newcommand{\Q} {\mathbb{Q}}
\newcommand{\R} {\mathbb{R}}
\newcommand{\Rb}{\overline{\mathbb{R}}}
\newcommand{\Rp}{\R_+}
\newcommand{\Rm}{\R_-}
\newcommand{\K} {\mathbb{K}}
\newcommand{\Ket} {\mathbb{K^{\ast}}}
\newcommand{\Cx}{\mathbb{C}}

%Ensembles, tribus
\newcommand{\A}{\mathbb{A}}
\renewcommand{\P}{\mathcal{P}}
\newcommand{\T}{\mathcal{T}}
\newcommand{\F}{\mathcal{F}}
\newcommand{\C} {\mathcal{C}}
\newcommand{\ouv}{\mathcal{O}}
\newcommand{\B}{\mathcal{B}}
\newcommand{\M}{\mathcal{M}}
\newcommand{\etag}{\mathcal{E}}
\newcommand{\etagOp}{\etag^{+}(\Omega)}
\newcommand{\etagO}{\etag(\Omega)}
%%Lp avant quotientage
\newcommand{\Lic}{\mathcal{L}^1}
\newcommand{\Lpc}{\mathcal{L}^p}
\newcommand{\Linfc}{\mathcal{L}^{\infty}}
\newcommand{\Lifull}{\Lic(\Omega, \B, m)}
\newcommand{\Lpfull}{\Lpc(\Omega, \B, m)}
\newcommand{\Linffull}{\Linfc(\Omega, \B, m)}
%%Lp après quotientage
\newcommand{\Li}{L^1}
\newcommand{\Ld}{L^2}
\newcommand{\Lp}{L^p}
\newcommand{\Linf}{L^{\infty}}
\newcommand{\Listd}{\Li(\Omega, m)} 
\newcommand{\Ldstd}{\Ld(\Omega, m)} 
\newcommand{\Lpstd}{\Lp(\Omega, m)} 

%Quelques opérateurs
\DeclareMathOperator{\codim}{codim}
\DeclareMathOperator{\rg}{rg}
\DeclareMathOperator{\tr}{tr}
\DeclareMathOperator{\vect}{Vect}
\DeclareMathOperator{\ima}{Im}
\DeclareMathOperator{\id}{Id}
\DeclareMathOperator{\sign}{sign}
\DeclareMathOperator{\ext}{ext}
\newcommand{\ind}{\mathbbm{1}}
\newcommand*{\spr}[2]{\left(\,#1\,\middle|\,#2\,\right)}
\newcommand{\interior}[1]{%
  {\kern0pt#1}^{\mathrm{o}}%
}
\DeclareMathOperator{\card}{card}
\DeclareMathOperator{\ppcm}{ppcm}

%Commandes de pure flemme
\newcommand{\veps}{\varepsilon}
\newcommand{\intab}{\int_{a}^{b}}
\newcommand{\intR}{\int_{-\infty}^{\infty}}
\newcommand{\intinf}{\int_{- \infty}^{+ \infty}}
\newcommand{\equi}{\Leftrightarrow}
\newcommand{\Kev}{$\K$-espace vectoriel }
\newcommand{\Kevs}{$\K$-espaces vectoriels }
\newcommand{\Rev}{$\R$-espace vectoriel }
\newcommand{\Revs}{$\R$-espaces vectoriels }
\newcommand{\Cev}{$\Cx$-espace vectoriel }
\newcommand{\Cevs}{$\Cx$-espaces vectoriels }
\newcommand{\evn}{(E, \norm{.})}
\newcommand{\interf}[2]{[#1,#2]}
\newcommand{\limb}{\lim\limits_}
\newcommand{\lebn}{\lambda_n}
\newcommand{\pp}{\text{~presque partout}}
\newcommand{\derivp}[2]{\frac{\partial #1}{\partial #2}}
\newcommand{\famiu}{(u_i)_{i \in I}}
\newcommand{\sev}{\text{~s.e.v.}}
\newcommand{\Mnp}{M_{n,p}(\K)}
\newcommand{\Mn}{M_{n}(\K)}
\newcommand{\tq}{\text{~t.q.~}}
\newcommand{\mq}{~montrer~que~}
\newcommand{\et}{\quad \text{et} \quad}
\newcommand{\ou}{\text{~ou~}}
\newcommand{\Kx}{\K[X]}


%Commandes de gars chelou
\renewcommand{\subset}{\subseteq}

%Commande restriction, ty duckduckgo
\def\restr#1#2{\mathchoice
              {\setbox1\hbox{${\displaystyle #1}_{\scriptstyle #2}$}
              \restraux{#1}{#2}}
              {\setbox1\hbox{${\textstyle #1}_{\scriptstyle #2}$}
              \restraux{#1}{#2}}
              {\setbox1\hbox{${\scriptstyle #1}_{\scriptscriptstyle #2}$}
              \restraux{#1}{#2}}
              {\setbox1\hbox{${\scriptscriptstyle #1}_{\scriptscriptstyle #2}$}
              \restraux{#1}{#2}}}
\def\restraux#1#2{{#1\,\smash{\vrule height .8\ht1 depth .85\dp1}}_{\,#2}} 

%Quelques normes
\DeclarePairedDelimiter\abs{\lvert}{\rvert}
\DeclarePairedDelimiter\labs{\bigl\lvert}{\bigr\rvert}
\DeclarePairedDelimiter\norm{\lVert}{\rVert}
\DeclarePairedDelimiterXPP\normi[1]{}\lVert\rVert{_1}{\ifblank{#1}{\:\cdot\:}{#1}}
\DeclarePairedDelimiterXPP\normd[1]{}\lVert\rVert{_2}{\ifblank{#1}{\:\cdot\:}{#1}}
\DeclarePairedDelimiterXPP\normp[1]{}\lVert\rVert{_p}{\ifblank{#1}{\:\cdot\:}{#1}}
\DeclarePairedDelimiterXPP\normq[1]{}\lVert\rVert{_q}{\ifblank{#1}{\:\cdot\:}{#1}}
\DeclarePairedDelimiterXPP\norminf[1]{}\lVert\rVert{_\infty}{\ifblank{#1}{\:\cdot\:}{#1}}

%Bracket en angle
\DeclarePairedDelimiter\angl{\langle}{\rangle}

%Set, parenthèses
\newcommand{\set}[1]{\{ #1 \}}
\newcommand{\bigset}[1]{\bigl\{ #1 \bigr\}}
\newcommand{\Bigset}[1]{\Bigl\{ #1 \Bigr\}}
\newcommand{\bigpar}[1]{\bigl( #1 \bigr)}
\newcommand{\Bigpar}[1]{\Bigl( #1 \Bigr)}

%Opitmisation des opérateurs de comparaison
\DeclarePairedDelimiter\tio{o (}{)}
\DeclarePairedDelimiter\gro{O (}{)}

\newcommand{\Gln}{\text{GL}(n,\K)}
\newcommand{\trans}[1]{\prescript{t}{}{#1}}

\pagestyle{fancy}

\fancyhead[C]{}
\fancyhead[R]{}

\begin{document}

\underline{Source :} Rombaldi - Algèbre - page 690 (LU - Cholesky) et 478 (CN, Sylvester)

\underline{Leçons :}
\begin{itemize}
\item  157 : Matrices symétriques réelles, matrices hermitiennes.
\item  170 : Formes quadratiques sur un espace vectoriel de dimension finie. Orthogonalité. Applications.
\item  171 : Formes quadratiques réelles. Coniques. Exemples et applications.
\item  162 : Systèmes d’équations linéaires ; opérations élémentaires, aspects algorithmiques et conséquences théoriques.
\end{itemize}

\begin{lem}
\label{lem:1}
Si elle existe, l'inverse d'une matrice triangulaire inférieure (resp. triangulaire supérieure) est triangulaire inférieure
(resp. triangulaire supérieure).
En particulier, si cette matrice est de diagonale unité, alors son inverse l'est aussi.
\end{lem}

\begin{proof}
Soit $A = (a_{i,j}) \in \Gln$ une matrice triangulaire inférieure. Notons $B = (b_{i,j}) \in \Gln$ son inverse.
Pour $(i,j) \in \set{1, \dots, n}^2$ donné, on a :
\[ \sum_{k=1}^n a_{i,k} b_{k,j} = \sum_{k=1}^i a_{i,k} b_{k,j} = \delta_{i,j} . \]
Pour $i=1$, cette relation donne directement $b_{1,2} = \dots = b_{1, n} = 0$. \\
Pour $i=2$, tenant compte du résultat précédent, on obtient de même $b_{2,3} = \dots = b_{2,n} = 0$. \\
En parcourant ainsi les valeurs de $i$ jusque $n$ on déduit que $B$ est triangulaire inférieure (la démonstration est identique
dans le cas triangulaire supérieure).
Si de plus $\forall i \in \set{1 \dots, n}, a_{i,i} = 1$ on déduit par calcul direct que tous les $b_{i,i}$ sont aussi égaux à 1.
\end{proof}

\begin{defn}[Sous-matrices principales]
Les sous-matrices principales de $A=(a_{i,j}) \in \Mn$ sont les matrices $(a_{i,j})_{1 \leq, i,j \leq k}$
et les déterminants principaux sont les déterminants associés.
\end{defn}

\begin{thm}[Décomposition LU]
Une matrice $A \in \Gln$ admet une décomposition $A = LU$ avec $L$ une matrice triangulaire inférieure de diagonale unité
et $U$ une matrice triangulaire supérieure si, et seulement si, tous les déterminants principaux de $A$ sont non nuls.
De plus, si une telle décomposition existe, elle est unique.
\end{thm}

\begin{proof}
On démontre d'abord l'équivalence par récurrence sur $n$. Le résultat est trivial pour $n=1$. 
Supposons qu'il soit vrai à un rang $n \geq 1$ quelconque et démontrons qu'il reste vrai au rang $n+1$. \\
Soit $A \in \Gln$ dont tous les déterminants principaux sont non nuls, on écrit :
\[ A = \begin{pmatrix} A_1 & B_1 \\ C_1 & a_{n,n} \end{pmatrix} , \]
avec $A_1 \in \text{GL}(n-1,\K)$ (par hypothèse sur $A$), $B_1 \in M_{1, n-1}(\K), C_1 \in M_{n-1,1}(\K)$ et $a_{n,n} \in \K$.
Par hypothèse de récurrence, on peut écrire $A_1 = L_1 U_1$ avec $L_1$ et $U_1$ telles que décrites dans l'énoncé.
Posons :
\[ L = \begin{pmatrix} L_1 & (0) \\ D_1 & 1 \end{pmatrix} \et U = \begin{pmatrix} U_1 & E_1 \\ (0) & \alpha \end{pmatrix} \]
avec $D_1 \in M_{1,n-1}(\K), E_1 \in  M_{n-1,1}(\K)$ et $\alpha \in \K$ à déterminer. 
L'équation $A=LU$ donne :
\[ \left \{ \begin{array}{c}
L_1 U_1 = A_1 \\
L_1 E_1 = B_1 \\
D_1 U_1  = C_1 \\
D_1 E_1 + \alpha = a_{n,n} 
 \end{array} \right . \]
Comme $L_1, U_1 \in \text{GL}(n-1,\K)$ ($\det(AB)=\det(A)\det(B)$) les matrices $E_1$ et $D_1$ existent et sont uniquement déterminées.
Par suite, il en est de même de $\alpha$. \\

Réciproquement, soit $A \in \Gln$ décomposable sous la forme $A=LU$ avec $L$ et $U$ telles qu'énoncées.
En reprenant la décomposition par blocs de $L$ et $U$ réalisée pour la condition suffisante, on peut alors écrire :
\[ A = \begin{pmatrix} A_1 & B_1 \\ C_1 & a_{n,n} \end{pmatrix} , \]
avec $A_1 = L_1 U_1$ satisfaisant l'hypothèse de récurrence, donc dont tous les déterminants principaux sont non nuls.
Ainsi, tous les déterminants principaux de $A$ sont inversibles ce qui démontre la condition nécessaire et, par suite,
l'hérédité du résultat. \newline

On démontre maintenant l'unicité. On écrit $A = L_1 U_1 = L_2 U_2$ où toutes ces matrices sont inversibles, on a alors
$L_1 L_2^{-1} = U_1^{-1} U_2$ avec la matrice du membre de gauche triangulaire inférieure et la matrice du membre de droite 
triangulaire supérieure (cf. lemme $\ref{lem:1}$).
Ces deux matrices étant égales, on déduit qu'elles sont diagonales de diagonale unitaire (toujours d'après le lemme $\ref{lem:1}$),
donc égales à $I_n$ ce qui démontre l'unicité.
\end{proof}

\begin{lem}
\label{lem:3}
Soit $A$ la matrice d'une forme quadratique $q$ dans une base quelconque. 
Si la forme $q$ est définie positive alors tous les déterminants principaux de $A$ sont strictement positifs.
\end{lem}

\begin{proof}
Notons $B=(e_1, \dots, e_n)$ la base telle que $A=M_B(q)$. 
Les sous-matrices principales de $A$ sont les matrices des restrictions de la forme quadratique $q$ aux sous-espaces vectoriels
engendrés par $\vect(e_1, \dots, e_k)$, pour $k \in \set{1, \dots, n}$, en particulier ces matrices sont toutes de déterminant 
strictement positif, ce qui démontre le résultat.
\end{proof}

\begin{rmq}
Une matrice définie positive est une matrice d'une forme quadratique définie positive et vice-versa.
\end{rmq}

\begin{thm}[Cholesky]
Une matrice $A$ à coefficients réels est symétrique et définie positive si, et seulement si, il existe une matrice $B$ triangulaire inférieure
et inversible telle que $A =B\trans{B}$. De plus, une telle décomposition est unique si on impose la positivité des coefficients diagonaux
de la matrice $B$. 
\end{thm}

\begin{proof}
Si $A$ est symétrique et définie positive alors, selon le lemme $\ref{lem:2}$ et le théorème précédent, $A$ admet une décomposition
$LU$ et de plus comme $\det{U} = \det{A} \neq 0$ on peut écrire $U = D U'$ avec $U'$ une matrice triangulaire supérieure 
de diagonale unité. On a ainsi $A=LDU'$ et comme $A=\trans{A}=U'^{T}D\trans{L}$, l'unicité de la décomposition $LU$ impose que 
$\trans{U'} = L$, ce qui permet d'écrire $A=LD\trans{L}$. 
Ainsi, la matrice $A$ étant définie positive, les coefficients $(d_i)$ de la matrice $D$ sont strictement positifs (on retrouve indirectement
qu'en dimension finie, il existe toujours une base $q$-orthogonale), 
ainsi on peut écrire $D = D'^2$ avec $D'$ une matrice diagonale dont les coefficients sont égaux à 
$\pm \sqrt{d_i}$ avec $d_i$ les coefficients de $D$. Par suite, $A=LD'D'\trans{L} = (LD')\trans{(LD')}$, ce qui est la forme attendue. 
De plus, la décomposition est bien unique si on impose la positivité des coefficients de la matrice $D'$. \\

Réciproquement, si $A=B \trans{B}$ avec $B$ inversible, alors pour toute matrice colonne $x$ non nulle, on a :
\[ \trans{x} A x = \trans{x} B \trans{B} x = \angl{\trans{B} x, \trans{B} x} > 0, \]
car $B$ est inversible (et donc $\trans{B} x$ non nul).
\end{proof}


\end{document}